% ============================================================
% 4  Near the transition: keep more terms -> Lifshitz theory
% ============================================================
\begin{frame}[t]
\ft{Near the transition: keep more terms}\par
\vspace{0.3cm}
\centering
% energy of a uniform magnetization m on the three sides of the transition (Landau picture)
\begin{tikzpicture}[font=\footnotesize,
  ax/.style={-{Stealth[length=3pt]},gray},
  curve/.style={line width=1.6pt,domain=-1.25:1.25,samples=60}]
 \foreach \xo/\c/\a/\b/\top/\bot/\st in
   {0/cLL/0.9/0.25/{small $\Delta V$}/{cost $\propto m^2/K_s$}/1,
    3.75/cQCP/0/0.8/{$\Delta V_c$}/{only $m^4$ left}/2,
    7.5/cFM/-1.0/0.6/{large $\Delta V$}/{$\langle m\rangle\neq0$}/3} {
   \onslide<\st->{%
   \begin{scope}[xshift=\xo cm]
   \fill[\c!10,rounded corners=4pt] (-1.55,-0.75) rectangle (1.55,1.45);
   \draw[ax] (-1.4,0)--(1.45,0) node[right,font=\scriptsize]{$m$};
   \draw[ax] (0,-0.6)--(0,1.3) node[left,font=\scriptsize]{$E$};
   \begin{scope} \clip (-1.55,-0.75) rectangle (1.55,1.25);
   \draw[\c,curve] plot (\x,{\a*\x*\x+\b*\x*\x*\x*\x}); \end{scope}
   \node[\c,font=\small\bfseries] at (0,1.7) {\top};
   \node at (0,-1.0) {\bot};
   \end{scope}}}
\end{tikzpicture}\par
\vspace{0.3cm}
\onslide<4->{%
{\small$\displaystyle
\mathcal L = i\,\partial_x\theta\,\partial_\tau\phi + \tfrac12(\partial_x\theta)^2
+ \underbrace{\colorbox{cLL!20}{$\mystrut\tfrac{\lambda_2}{2}(\partial_x\phi)^2$}}_{\text{\scriptsize\color{cLL} Luttinger liquid}}
+ \underbrace{\colorbox{cQCP!15}{$\mystrut\tfrac12(\partial_x^2\phi)^2 + \tfrac{\lambda_4}{4}(\partial_x\phi)^4$}}_{\text{\scriptsize\color{cQCP} next terms}}
+ \dots$\par}
\vspace{0.05cm}
{\scriptsize $\phi=\phi_s/\sqrt{2\pi}$,\ \ $\partial_x\phi\propto m$,\qquad $\lambda_2=(2\pi/K_s)^2\to0$ at $\Delta V_c$\par}}
\vspace{0.2cm}
\onslide<5->{%
{\small At $\Delta V_c$: \colorbox{orchid!25}{\mystrut Lifshitz theory}, $\omega\propto k^2$ ($z\approx2$): not slide 1's $\omega^2+k^2$\par}
\vspace{0.05cm}
{\small $\lambda_4$ relevant: \colorbox{pumpkin!40}{\mystrut interacting}, non-conformal critical point\par}
\vspace{0.1cm}
{\tiny\color{citegray} K.~Yang, PRL 93, 066401 (2004);\; A.~Nahum, arXiv:2506.21540 (2025);\; R.~Flores-Calder\'on, M.~Z\"undel, arXiv:2511.09097 (2025)\par}}
\onslide<5->{\tk{At the transition the Luttinger term vanishes: the next terms take over}}
\end{frame}
% Notes: Steps 1-3: The energy of a uniform magnetization m, Landau style. Small Delta V (Luttinger liquid): a parabola, the m^2 term (its coefficient ~1/K_s) is all we need: that is the Luttinger liquid, the first term of a Taylor expansion. At Delta V_c the m^2 coefficient vanishes (K_s -> infinity): the curve is flat at the bottom, only m^4 holds it. Beyond: a double well, the ferromagnet with <m> != 0.
% Notes: Step 4: In fields (paper's normalization, phi = phi_s / sqrt(2 pi), d_x phi ~ m): the Luttinger liquid keeps only lambda_2 (d_x phi)^2 with lambda_2 = (2 pi / K_s)^2. At the transition we must keep the next terms: (d_x^2 phi)^2 (gradient cost of the magnetization) and (d_x phi)^4 (the m^4 of the Landau curve).
% Notes: Step 5: This is the Lifshitz theory (K. Yang for this model; Nahum). Its quadratic part gives omega ~ k^2: z = 2 instead of 1, so the omega^2 + k^2 of slide 1 no longer holds: Mermin-Wagner's log divergence is gone. The quartic term is relevant: an interacting, non-conformal critical point; z = 2 - 2 eps^2/243 ~ 1.99 (Flores-Calderon & Zundel), consistent with numerics (Chepiga & Essler).
